\documentclass[11pt]{article}
\usepackage[margin=1in]{geometry}
\usepackage{amsmath}
\usepackage{booktabs}

\title{Step 3 Validation Protocol: Dynamic Dimerisation/Phonon Coordinate}
\author{}
\date{}

\begin{document}
\maketitle

\section{Purpose of the test}
Step 2 showed that a scalar effective mixing,
\begin{equation}
V_{BD}^{\mathrm{eff}} =
V_0 + \lambda_\delta \delta + \lambda_C C_1,
\end{equation}
can detect enhanced bright-dark mixing but cannot identify whether the
enhancement comes from spin-Peierls dimerisation or local spin correlations.
The purpose of Step 3 is to add a dynamic dimerisation/phonon coordinate and
test whether this creates a spectroscopic signature that a static $C_1$ scalar
control cannot reproduce.

\section{Physical model}
The orbital sector is the same minimal bright/dark model used in Steps 1 and
2.  The Hilbert space contains a ground state $|g\rangle$, a dark excited state
$|D\rangle$, a bright excited state $|B\rangle$, and a truncated harmonic
coordinate with annihilation operator $a$.  The phonon displacement is
\begin{equation}
Q = a + a^\dagger.
\end{equation}
The Step 3 mixing Hamiltonian is
\begin{equation}
H_{\mathrm{mix}} =
V_{BD}^{\mathrm{static}} L_{BD}
+ g_Q Q L_{BD},
\end{equation}
where
\begin{equation}
L_{BD} = |D\rangle\langle B| + |B\rangle\langle D|
\end{equation}
and
\begin{equation}
V_{BD}^{\mathrm{static}} =
V_0 + \lambda_\delta \delta + \lambda_C C_1.
\end{equation}
The new ingredient is not another scalar amplitude; it is the operator-valued
term $g_Q Q L_{BD}$.

\section{Expected results}
At matched $V_{BD}^{\mathrm{static}}$, the static dimerisation and static
spin-correlation controls should remain equivalent, reproducing the degeneracy
found in Step 2.  The dynamic dimerisation/phonon case is expected to differ
because it can mix phonon manifolds.  Possible signatures include phonon
sidebands, modified pathway amplitudes, altered peak positions, or changes in
rephasing/unrephasing balance.

\section{Numerical procedure}
Run one scenario at a time by setting \texttt{STEP3\_ACTIVE\_CASE} or by
editing \texttt{active\_case\_key} in the notebook.  The default scenarios are:
\begin{center}
\begin{tabular}{lll}
\toprule
Scenario & Purpose & Dynamic phonon term \\
\midrule
\texttt{static\_spin\_correlation} & scalar $C_1$ control & off \\
\texttt{static\_dimerisation} & scalar $\delta$ control & off \\
\texttt{dynamic\_dimerisation\_phonon} & dynamic lattice control & on \\
\bottomrule
\end{tabular}
\end{center}
The third-order rephasing 1Q spectrum is generated using the same NQ protocol
as in Steps 1 and 2.  The frequency window is widened to leave room for phonon
sideband structure.

\section{Acceptance criteria}
The first criterion is that the two static controls remain identical or nearly
identical at matched $V_{BD}^{\mathrm{static}}$.  This confirms that the Step 2
degeneracy is reproduced.

The second criterion is that the dynamic dimerisation/phonon case differs from
the static controls in at least one observable: cross-window area,
cross/diagonal ratio, pathway-resolved amplitude, peak position, linewidth,
or sideband intensity.  If the dynamic case remains indistinguishable, then the
chosen $g_Q$, phonon energy, Hilbert-space truncation, or observable is not
sufficient to isolate lattice dynamics.

\section{Requested figures and data}
For each run, save:
\begin{itemize}
\item the rephasing real spectrum PDF;
\item the pathway-resolved complex response arrays;
\item the total rephasing, unrephasing, and 1Q components;
\item text and JSON summaries containing $V_{BD}^{\mathrm{static}}$, $g_Q$,
$\omega_Q$, $\delta$, and $C_1$.
\end{itemize}
Phase 2 analysis should compare finite-window integrals, cross/diagonal ratios,
pairwise spectrum differences, and possible sideband windows offset by
$\omega_Q$.

\section{Interpretation of possible failures}
If the dynamic phonon case differs only by an overall amplitude, the observable
may still be insensitive to true lattice dynamics.  If the static controls do
not match at equal $V_{BD}^{\mathrm{static}}$, the model implementation should
be checked before interpreting the dynamic case.  If sidebands are absent, the
phonon coupling may be too weak, the frequency window too narrow, or the
truncated phonon basis too small.

\section{Numerical results}
The three Step 3 scenarios were generated at $\eta=0.001$ with a
$140\times140$ frequency grid, $n_{\mathrm{bosons}}=3$, and
$\omega_Q=0.035$ eV.  The finite-window half-width used for the aggregate
metrics is $0.025$ eV.  The two static controls were run at matched
$V_{BD}^{\mathrm{static}}=0.02$.

\begin{center}
\begin{tabular}{lrrrr}
\toprule
Scenario & $V_{BD}^{\mathrm{static}}$ & $g_Q$ &
$A_{\mathrm{cross}}$ & $A_{\mathrm{cross}}/A_{\mathrm{diag}}$ \\
\midrule
Static spin correlation & 0.0200 & 0 & 14.160315 & 0.142697 \\
Static dimerisation & 0.0200 & 0 & 14.160315 & 0.142697 \\
Dynamic dimerisation phonon & 0.0200 & 0.0100 & 11.761513 & 0.149235 \\
\bottomrule
\end{tabular}
\end{center}

The static dimerisation and static spin-correlation spectra are identical to
numerical precision.  Their relative maximum full-spectrum difference is
$1.10\times10^{-17}$.  This reproduces the Step 2 scalar degeneracy inside the
larger phonon Hilbert space.

The dynamic dimerisation/phonon spectrum differs strongly from the static
dimerisation control at the same $V_{BD}^{\mathrm{static}}$.  The relative
maximum full-spectrum difference is $5.46\times10^{-1}$.  The main cross-window
area decreases from $14.160315$ to $11.761513$, while the cross/diagonal ratio
increases from $0.142697$ to $0.149235$.

The sideband-window metric is qualitative because the finite windows can include
tails from nearby resonances.  Nevertheless, it also changes between the static
and dynamic cases: the sideband/main-cross ratio changes from $2.21297$ for the
static controls to $2.41282$ for the dynamic phonon case.

\section{Validation outcome}
The first criterion passes.  The two static controls are identical at matched
$V_{BD}^{\mathrm{static}}$ within relative maximum difference
$1.10\times10^{-17}$.  This confirms that simply embedding the scalar model in
a larger phonon Hilbert space does not by itself break the degeneracy between
$\delta$ and $C_1$.

The second criterion also passes.  Adding the dynamic term $g_Q Q L_{BD}$
produces a distinct nonlinear spectrum relative to the static dimerisation
control, with relative maximum full-spectrum difference $5.46\times10^{-1}$.
Therefore an explicit dimerisation/phonon coordinate provides a viable route to
separate lattice-driven spin-Peierls physics from a purely static local
spin-correlation control.

\section{Conclusions and next step}
Step 3 resolves the main limitation of Step 2 at the model level.  A scalar
renormalisation of bright-dark mixing cannot identify the microscopic origin of
the signal, but an operator-valued phonon coordinate can generate a distinct
spectral response.  This does not yet prove that a specific feature in CuGeO3 is
phonon-driven, but it establishes a workable modelling criterion: lattice
physics must enter through a dynamical or sector-specific degree of freedom,
not only through a scalar $V_{BD}^{\mathrm{eff}}$.

The next useful calculation is not to increase grid convergence immediately.
Instead, the next step should scan the dynamical parameters one at a time:
\begin{itemize}
\item $g_Q$ to determine the onset and approximate scaling of the dynamic
phonon signature;
\item $\omega_Q$ to check whether sideband displacement follows the phonon
energy;
\item rephasing versus unrephasing totals to see whether the dynamic coordinate
changes pathway balance;
\item real, imaginary, and absolute spectra to identify the most robust
observable.
\end{itemize}
After that, the model should be connected to temperature and field trends by
using $\delta(T,B)$ for the dynamic lattice channel and $C_1(T,B)$ for the
static spin-correlation control.

\end{document}
